%% Draft-only owner placeholder (ruling R15, 2026-09-09). Prints ONLY when the
%% build defines \DRAFTTODO (e.g. pdflatex "\def\DRAFTTODO{1}%% Working draft. Structure is generated; prose is authored here.
%% Build from the repository root with: bash scripts/build_paper.sh <this paper dir>
\documentclass[11pt]{article}

\usepackage[utf8]{inputenc}
\usepackage[T1]{fontenc}
\usepackage{amsmath,amssymb}
\usepackage{booktabs}
\usepackage{graphicx}
\usepackage{tikz}
\usepackage[margin=1in]{geometry}
% For every included figure, pdfTeX records the absolute path it was read
% from in a /PTEX.FileName object.  Those objects are not rendered, so a
% text dump of the PDF does not show them, but they travel with the file
% and spell out the build machine's home directory, user name and worktree
% layout.  Suppressing them costs nothing: they exist for debugging.
\pdfsuppressptexinfo=-1
\usepackage[hidelinks]{hyperref}

\graphicspath{{../figs/out/}}

\renewcommand{\topfraction}{0.9}
\renewcommand{\bottomfraction}{0.7}
\renewcommand{\textfraction}{0.08}
\renewcommand{\floatpagefraction}{0.8}
\setcounter{topnumber}{2}
\setcounter{totalnumber}{3}

%% Numbers and tables are emitted by the figure code from hash-bound evidence.
% Generated by the figure script from hash-bound evidence. Do not hand-edit.
\newcommand{\ToyN}{10}
\newcommand{\ToySeed}{0}
\newcommand{\ToySteps}{50}
\newcommand{\ToyLongSteps}{500}
\newcommand{\BudgetMultiple}{ten}
\newcommand{\BudgetFold}{tenfold}
\newcommand{\ToyEvalSteps}{30}
\newcommand{\ToyRuns}{1}
\newcommand{\ToyDegraded}{20.75}
\newcommand{\ToyDegradedSSIM}{0.460}
\newcommand{\ToyConv}{14.72}
\newcommand{\ToyConvSSIM}{0.190}
\newcommand{\ToyLong}{14.05}
\newcommand{\ToyLongSSIM}{0.242}
\newcommand{\ToyTrans}{6.57}
\newcommand{\ToyTransSSIM}{0.005}
\newcommand{\ToyConvShortfall}{6.03}
\newcommand{\ToyTransShortfall}{14.18}
\newcommand{\ToyConvParams}{16,058,179}
\newcommand{\ToyTransParams}{17,961,024}
\newcommand{\ToyParamExcess}{11.8}
\newcommand{\ToyConvSeconds}{47.46}
\newcommand{\ToyTransSeconds}{1.58}
\newcommand{\ToySpeedRatio}{29.9}
\newcommand{\ToyConvLoss}{0.0810}
\newcommand{\ToyTransLoss}{0.9912}
\newcommand{\ToyTransLastTen}{0.9880}
\newcommand{\ToyLongLoss}{0.0338}
\newcommand{\ToyLongSeconds}{478.17}
\newcommand{\ToyLossDrop}{2.4}
\newcommand{\ToyLongRegression}{0.67}
\newcommand{\SubsetN}{100}
\newcommand{\SubsetCameras}{2}
\newcommand{\MethodPSNR}{21.88}
\newcommand{\MethodSSIM}{0.614}
\newcommand{\MethodLow}{21.41}
\newcommand{\MethodHigh}{22.35}
\newcommand{\BicubicPSNR}{25.48}
\newcommand{\BicubicSSIM}{0.709}
\newcommand{\NearestPSNR}{24.65}
\newcommand{\NearestSSIM}{0.674}
\newcommand{\MethodYPSNR}{23.44}
\newcommand{\BicubicYPSNR}{27.11}
\newcommand{\GapBicubic}{$-3.60$}
\newcommand{\GapBicubicAbs}{3.60}
\newcommand{\GapNearest}{$-2.77$}
\newcommand{\GapSSIM}{$-0.095$}
\newcommand{\GapLower}{$-4.06$}
\newcommand{\GapUpper}{$-3.19$}
\newcommand{\GapRecordedLower}{$-4.04$}
\newcommand{\GapRecordedUpper}{$-3.12$}
\newcommand{\GapBest}{$+4.02$}
\newcommand{\GapWorst}{$-11.05$}
\newcommand{\WinsBicubic}{3}
\newcommand{\LossesBicubic}{97}
\newcommand{\WinsNearest}{9}
\newcommand{\LossesNearest}{91}
\newcommand{\GapSignP}{$2.6 \times 10^{-25}$}
\newcommand{\GapWilcoxP}{$5.4 \times 10^{-17}$}
\newcommand{\GapTLower}{$-4.06$}
\newcommand{\GapTUpper}{$-3.14$}
\newcommand{\GapNearestSignP}{$3.3 \times 10^{-18}$}
\newcommand{\NoiseMax}{0.098}
\newcommand{\NoiseMean}{0.011}
\newcommand{\NoiseShift}{$+0.0028$}
\newcommand{\BicubicRecordState}{the same bytes at two paths}
\newcommand{\CensusTotal}{206}
\newcommand{\CensusRan}{13}
\newcommand{\CensusSkipped}{193}
\newcommand{\CensusSkipShare}{93.7}
\newcommand{\CensusCameras}{9}
\newcommand{\AbsentGroups}{7}
\newcommand{\BudgetPixels}{393,216}
\newcommand{\BudgetLong}{768}
\newcommand{\BudgetVram}{32}
\newcommand{\SkippedSide}{2048}
\newcommand{\SkippedPixels}{4,194,304}
\newcommand{\BudgetExcess}{10.7}
\newcommand{\OrphanStems}{13}
\newcommand{\CensusPathState}{no longer matches the path it names}
\newcommand{\EvalRun}{100}
\newcommand{\EvalSkipped}{0}
\newcommand{\EvalReferences}{100}
\newcommand{\PrimaryTier}{blocked}
\newcommand{\SotaTier}{blocked}
\newcommand{\LpipsWins}{93}
\newcommand{\LpipsLosses}{7}
\newcommand{\LpipsTies}{0}
\newcommand{\LpipsSignP}{$2.7 \times 10^{-20}$}
\newcommand{\LpipsDit}{0.317}
\newcommand{\LpipsBicubic}{0.456}
\newcommand{\LpipsDelta}{0.139}
\newcommand{\LpipsDeltaLo}{0.118}
\newcommand{\LpipsDeltaHi}{0.160}
\newcommand{\LpipsCanonWins}{46}
\newcommand{\LpipsNikonWins}{47}
\newcommand{\QuadLossWin}{89}
\newcommand{\QuadLossLoss}{7}
\newcommand{\QuadWinWin}{4}
\newcommand{\QuadWinLoss}{0}
\newcommand{\AnchorDeltaDb}{0.07}
\newcommand{\AnchorLocalPsnr}{21.95}
\newcommand{\AnchorLosses}{96}
\newcommand{\PsnrLossesLocal}{96}
\newcommand{\PsnrWinsLocal}{4}
\newcommand{\GapBicubicLocal}{$-3.52$}
\newcommand{\LpipsSeed}{20260909}
\newcommand{\LpipsSteps}{40}
\newcommand{\LpipsVggDit}{0.416}
\newcommand{\LpipsVggBicubic}{0.407}
\newcommand{\LpipsVggDelta}{$-0.009$}
\newcommand{\LpipsVggDeltaLo}{$-0.021$}
\newcommand{\LpipsVggDeltaHi}{0.003}
\newcommand{\LpipsVggWins}{41}
\newcommand{\LpipsVggLosses}{59}
\newcommand{\LpipsVggTies}{0}
\newcommand{\LpipsVggSignP}{0.0886}
\newcommand{\LpipsVggCanonWins}{18}
\newcommand{\LpipsVggNikonWins}{23}
\newcommand{\QuadVggLossWin}{38}
\newcommand{\QuadVggLossLoss}{58}
\newcommand{\QuadVggWinWin}{3}
\newcommand{\QuadVggWinLoss}{1}
\newcommand{\InitCells}{9}
\newcommand{\InitSeeds}{3}
\newcommand{\InitSteps}{50}
\newcommand{\InitThreshold}{0.90}
\newcommand{\InitDefaultStall}{3}
\newcommand{\InitStandardStall}{3}
\newcommand{\InitZeroEscape}{3}
\newcommand{\InitZeroLoss}{0.42}
\newcommand{\InitDefaultLoss}{0.99}
\newcommand{\InitStandardLoss}{0.99}
\newcommand{\InitStallLoss}{0.99}
\newcommand{\InitZeroPsnr}{9.7}
\newcommand{\InitStallPsnr}{6.6}
\newcommand{\InitZeroGain}{3.2}
\newcommand{\IdentityPsnr}{20.75}
\newcommand{\NEvidence}{37}
\newcommand{\EvidenceBytes}{416,637}

%% Draft-only owner placeholder (ruling R15, 2026-09-09). Prints ONLY when the
%% build defines \DRAFTTODO (e.g. pdflatex "\def\DRAFTTODO{1}%% Working draft. Structure is generated; prose is authored here.
%% Build from the repository root with: bash scripts/build_paper.sh <this paper dir>
\documentclass[11pt]{article}

\usepackage[utf8]{inputenc}
\usepackage[T1]{fontenc}
\usepackage{amsmath,amssymb}
\usepackage{booktabs}
\usepackage{graphicx}
\usepackage{tikz}
\usepackage[margin=1in]{geometry}
% For every included figure, pdfTeX records the absolute path it was read
% from in a /PTEX.FileName object.  Those objects are not rendered, so a
% text dump of the PDF does not show them, but they travel with the file
% and spell out the build machine's home directory, user name and worktree
% layout.  Suppressing them costs nothing: they exist for debugging.
\pdfsuppressptexinfo=-1
\usepackage[hidelinks]{hyperref}

\graphicspath{{../figs/out/}}

\renewcommand{\topfraction}{0.9}
\renewcommand{\bottomfraction}{0.7}
\renewcommand{\textfraction}{0.08}
\renewcommand{\floatpagefraction}{0.8}
\setcounter{topnumber}{2}
\setcounter{totalnumber}{3}

%% Numbers and tables are emitted by the figure code from hash-bound evidence.
\input{generated_numbers}
\input{captions}

\title{Print the trivial arm: three scales at which a do-nothing baseline
decided what an image restoration comparison could conclude}
\author{Working draft --- author list not yet fixed}
\date{Working draft. Not submitted to any venue.}

\begin{document}
\maketitle

\begin{abstract}
\input{abstract}
\end{abstract}

\input{body}

\bibliographystyle{plain}
\bibliography{references}

\end{document}
");
%% every packet build leaves it undefined, so no unconfirmed declaration text
%% reaches a PDF. The owner replaces each use with the confirmed sentence.
\providecommand{\ownertodo}[1]{\ifdefined\DRAFTTODO[TODO-OWNER: #1]\fi}

%% Self-contained captions. No cross-references, no section numbers: each
%% caption must be readable on its own if the figure is lifted out.

\newcommand{\CapScales}{%
The three scales at which the cheapest arm decided what the comparison could
conclude, placed on one page before those scales are read separately. The
Toy-scale panel restates the toy run: the degraded input sits at
\ToyDegraded{} decibels and the best trained arm at \ToyConv{}, so both
architectures are below doing nothing. The Longer-budget panel restates the
longer budget: the objective falls and fidelity falls with it to \ToyLong{},
still below the untouched input. The Crop-128 panel restates the measured
subset: the released model is behind bicubic on \LossesBicubic{} of
\SubsetN{} crops, by \GapBicubicAbs{} decibels on average, on peak
signal-to-noise ratio. All three panels are fidelity readouts; the perceptual
endpoint on which the subset ranking reverses is drawn on its own figure, and
no panel adds an official ground-truth score.%
}

\newcommand{\CapToy}{%
Three trained arms, all below the input they were asked to repair. A toy-scale
restoration run trained a small convolutional network and a small transformer of
comparable size under a matched step budget on the same data, and scored both
against the same \ToyN{} held-out pairs at \ToyEvalSteps{} sampling steps. Points
are mean peak signal-to-noise ratio of the restored output against the reference
image; higher is closer to the reference. The vertical rule is the same quantity
for the degraded input passed through unchanged, which the same record scores as
an arm in its own right. No trained arm reaches it, so the run supports no
comparison between the two architectures: the shortfall of the best of them is
\ToyConvShortfall{} decibels. The transformer arm's training objective sat at
\ToyTransLoss{} at the last step, against a mean of \ToyTransLastTen{} over the
last ten steps and \ToyConvLoss{} for the convolutional arm at the same step
count: it did not fall, and a denoising objective that stays at unity is what a
network whose output does not depend on its input produces. The stall
reproduces in \InitDefaultStall{} of \InitSeeds{} seeds of the same construction
in a separate predeclared ablation, which also shows that the one
initialisation that zeroes the network's residual gates exactly is the one that
trains; the earlier attribution of the stall to zero-initialised gates is
withdrawn.%
}

\newcommand{\CapBudget}{%
The training budget raised \BudgetFold{}, read on both axes at once. The same
small convolutional restoration network was trained for \ToySteps{} and for
\ToyLongSteps{} steps on the same data at the same learning rate and seed, and
evaluated on the same \ToyN{} held-out pairs at \ToyEvalSteps{} sampling steps.
The two facets carry their own labels: one is the training objective at the last
step, the other the mean peak signal-to-noise ratio of the restored output
against the reference image, where the horizontal rule marks the same quantity
for the degraded input passed through unchanged. The objective falls by a factor
of \ToyLossDrop{}. The fidelity of the output falls with it, by
\ToyLongRegression{} decibels, and stays below the untouched input at both
budgets. Read alone, either facet would give the opposite impression of what the
extra budget bought. The longer run was
written outside the results tree and was never promoted.%
}

\newcommand{\CapSubset}{%
Every crop, the released model against an interpolation. Each point is one crop
of a \SubsetN{}-crop subset drawn from a real-world super-resolution validation
set. The horizontal position is the peak signal-to-noise ratio of bicubic
upsampling of the low-resolution input against the reference image; the vertical
position is the same quantity for a released diffusion restoration model given
the same input, scored by the same metric pass over the same crops. The diagonal
is equality. Points below it are crops on which the interpolation is closer to
the reference. Symbols distinguish the two camera groups the subset is drawn
from, in equal numbers. Peak signal-to-noise ratio measures distance from the
reference and not perceptual quality, so the panel reports which output is
closer to the reference on that measure and not which is better; on a learned
perceptual distance scored on the same crops the order reverses, and that
endpoint is drawn on its own figure. Filled and open symbols identify the two
camera groups; the diagonal is the pre-specified equality reference, not a
fitted trend.%
}

\newcommand{\CapPaired}{%
The paired differences, and the width the model varies by against itself. Panel
A plots one point per crop: the peak signal-to-noise ratio of a released
diffusion restoration model minus that of bicubic upsampling of the same
low-resolution input, over a \SubsetN{}-crop subset, ordered by that difference.
Negative values are crops on which the interpolation is closer to the reference.
The solid rule is zero and the dashed rule is the mean difference. The shaded
band is the largest per-crop shift between two inference passes of the same
model on the same crops, \NoiseMax{} decibels. The record does not show the
sampling seed varied between the passes, so two runs cannot be read as an
estimate of sampling variance; the band bounds only how far re-running the
recorded inference moved a crop. Panel B
gives the mean difference overall and within each camera group, with 95\%
percentile bootstrap intervals over the crops. Camera group is a recorded
property of each crop and was fixed before these differences were computed. The
zero rule is the equality point; the dashed rule is the observed mean, and the
interval in Panel B is a bootstrap summary over this measured subset rather
than uncertainty for the full census.%
}

\newcommand{\CapCeiling}{%
What a memory ceiling selected in an earlier census, and what it removed. That
census ran a released diffusion restoration model over the full set of
\CensusTotal{} inputs
under a fixed budget of \BudgetPixels{} output pixels and a maximum output side
of \BudgetLong{} pixels; inputs whose upsampled output would exceed it were
skipped before inference. Panel A gives, for that census and for the
later \SubsetN{}-crop published crop128 run, what became of every input the run
was given: dropped
before inference, run but left with no reference image to score against, or run
and scored. The census drop and the absence of references were not sampling
decisions; the later scored set is the published protocol, not the remnant of
the skip. Panel B breaks the dropped census inputs down by the recorded source
group
and shades the two groups that appear in the later crop128 set. Every dropped
input
was of a single larger size class, so what survived
the ceiling in the census was fixed by output dimensions rather than chosen.
Segment labels
are n/N and percentages: in Panel A, N is the number of inputs given to that
attempt, and in Panel B it is the total number dropped for memory. An
input that ran without a matched reference is counted as unscored rather than
treated as a zero-fidelity observation; zero-height outcomes remain represented
by the legend.%
}

\newcommand{\CapMetricConsistency}{%
Four fidelity metrics on one paired population. Each point is the mean
DiT--bicubic difference over the same \SubsetN{} crop128 pairs; the vertical
rule is equality and each horizontal whisker is a 95\% percentile bootstrap
interval over those paired differences. The facet labelled PSNR keeps RGB and
luma peak signal-to-noise ratio on their native decibel scale; the facet
labelled SSIM keeps RGB and luma structural similarity on its native unitless
scale. Negative
values mean that the model is farther from the reference under the displayed
fidelity measure. The facets deliberately use separate axes, so their
magnitudes are not compared across units; all four are distortion measures,
and the perceptual endpoint is drawn on its own figure.%
}

\newcommand{\CapCameraSensitivity}{%
Camera-stratified sensitivity of the paired RGB peak signal-to-noise ratio.
Each point is the mean DiT--bicubic difference within one recorded camera
stratum; the horizontal interval is the corresponding descriptive percentile
bootstrap over paired crops, and the label gives the number of DiT wins out of
the stratum's crops. The vertical rule is equality. Both strata are part of the
same 100-crop published crop128 protocol, so this panel does not provide
scene-level or full-census uncertainty and does not measure perceptual quality.%
}

\newcommand{\CapCropCameraMetrics}{%
Crop-level fidelity and scope diagnostics for the published crop128 subset.
Panel A shows the DiT--bicubic difference for every one of the \SubsetN{} paired
crops on RGB PSNR and RGB SSIM, with Canon and Nikon symbols retained. Panel B
gives descriptive 95\% bootstrap intervals for the same two metrics within
each camera stratum; these are intervals over the measured crops, not a
scene-level or full-census uncertainty. Panel C keeps the earlier full-input
census visible: the early \CensusTotal{}-input attempt yielded \CensusRan{} runs
and \CensusSkipped{} skips, whereas the separate crop128 protocol scored
\SubsetN{} crops with references. PSNR and SSIM measure fidelity to a
reference; this figure adds no perceptual endpoint, which is drawn on its own
figure. The metric rows
are recomputed from the bound DiT and bicubic per-crop score files; the bound
camera-stratum record supplies the recorded camera counts and means, and the
census is bound by the earlier evaluation and skip records.}

\newcommand{\CapToyTable}{%
The toy-scale arms as their own training records describe them. Parameter
counts, step counts, wall clock and the training objective at the last step are
read from the manifest each run wrote; fidelity is mean peak signal-to-noise
ratio of the restored output against the reference image over the same
\ToyN{} held-out pairs at \ToyEvalSteps{} sampling steps, and the first row is
the degraded input passed through unchanged, which has no training record.
The two arms trained at \ToySteps{} steps differ by \ToyParamExcess{} per cent in
parameter count, so the budget they share is a step budget and not a capacity
budget, and their wall clocks differ by a factor of \ToySpeedRatio{}, so it is
not a wall-clock budget either.%
}

\newcommand{\CapSubsetTable}{%
The subset comparison, all four recorded fidelity measures. Rows are the two
interpolations and the released diffusion restoration model, each scored over the
same \SubsetN{} crops against the same reference images by the same metric pass.
Peak signal-to-noise ratio and structural similarity are reported over the full
colour image and again over the luma channel with a four-pixel border removed.
The two final rows are paired differences, computed crop by crop and then
averaged, so a negative value means the model was further from the reference than
the interpolation on average. All four measures compare an output with its
reference and none of them measures perceptual quality; the learned perceptual
distance scored on the same crops, on which the order reverses, is reported in
its own figure rather than in this table because it was scored on a fresh
sample of the model rather than on the pixel set behind these rows.%
}

\newcommand{\CapPerceptualReversal}{%
The same crops on both axes of the perception--distortion plane. Each point is
one of the \SubsetN{} crops of the published RealSR crop128 protocol. The
horizontal position is the peak signal-to-noise ratio of a fresh sample of the
released diffusion restoration model, drawn at sampling seed \LpipsSeed{}, minus
that of bicubic upsampling of the same low-resolution input, so negative values
are crops on which the interpolation is closer to the reference; the model
loses on \PsnrLossesLocal{} of \SubsetN{}, and the fresh sample's mean sits
\AnchorDeltaDb{} decibels from the bound fidelity record under a predeclared
anchor rule. The vertical position is the learned perceptual distance (LPIPS,
AlexNet backbone, lower is closer) of bicubic minus that of the model on the
same crop and reference, so positive values are crops on which the model is
perceptually closer; the model wins on \LpipsWins{} of \SubsetN{}, with
\LpipsTies{} ties and an exact sign test of \LpipsSignP{}. The two rules are
equality on each axis and divide the crops into four quadrants, whose counts are
printed: on \QuadLossWin{} crops the interpolation is closer on fidelity and
the model is closer perceptually, and the two endpoints agree on only
\QuadLossLoss{} plus \QuadWinWin{} crops. Filled and open symbols are the two
camera groups. Neither axis is a verdict on the model; together they are the
trade-off, observed.%
}

\newcommand{\CapInitAblation}{%
The toy transformer's training objective under three initialisations at three
seeds each. Every cell is the same network, data, batch, learning rate and
\InitSteps{}-step budget as the bound toy run; only the initialisation and the
seed change. The three facets are the library's default construction, whose
gates are not zero and which is the arm behind the bound record; a standard
transformer scheme without any zeroing step, whose gates are not zero; and that
scheme plus exact zeroing of the adaptive-normalisation gates and final
projection. Each line joins the three points the record holds for one seed: the
mean objective over the first ten steps, the mean over the last ten, and the
last step. The solid rule is unity, the value a denoising objective takes when
the network's output does not depend on its input; the dashed rule is the
predeclared stall threshold of \InitThreshold{} on the last-ten mean. The
default stalls in \InitDefaultStall{} of \InitSeeds{} seeds and the standard
scheme in \InitStandardStall{} of \InitSeeds{}; the zero-gated scheme escapes
in \InitZeroEscape{} of \InitSeeds{}, to a last-ten mean of \InitZeroLoss{}
against \InitStallLoss{} for the cells that stalled. All nine cells restore to
below the degraded input at \IdentityPsnr{} decibels: \InitZeroPsnr{} for the
arm that trained against \InitStallPsnr{} for the arms that did not. The panel
supports a statement about this toy's initialisation and none about the
architecture.%
}

\newcommand{\CapCensusTable}{%
Which inputs a memory ceiling removed, by source group. An earlier attempt over
the full set of \CensusTotal{} inputs skipped every input whose upsampled output
would have exceeded \BudgetPixels{} pixels, which removed \CensusSkipped{} of
them. Groups are the recorded source labels, which correspond to the capture
device. The final column records whether that group appears in the
\SubsetN{}-crop subset measured later. Of the \CensusCameras{} groups the ceiling
touched, \AbsentGroups{} are absent from every measured result in this study.%
}


\title{Print the trivial arm: three scales at which a do-nothing baseline
decided what an image restoration comparison could conclude}
\author{Working draft --- author list not yet fixed}
\date{Working draft. Not submitted to any venue.}

\begin{document}
\maketitle

\begin{abstract}
Image restoration results are reported as comparisons between methods, and the
arm that does nothing is usually left out because its number is uninteresting.
This paper reports three comparisons from one project in which that number was
recorded anyway, and shows that in each case it decided what the comparison
could conclude. At toy scale, two small networks were trained under a
matched-step budget and compared with each other; both finished below the
degraded input they were asked to repair, at \ToyConv{} and \ToyTrans{} decibels
against \ToyDegraded{}, so no comparison between them was available, and the
project withdrew its own recorded conclusion once that was noticed. Raising the
training budget \BudgetFold{} on the better of the two cut the training objective by a
factor of \ToyLossDrop{} and moved restored fidelity \ToyLongRegression{}
decibels the wrong way. At subset scale, a released diffusion restoration model
was run to completion over \SubsetN{} crops with reference images present, and
lost to bicubic upsampling on \LossesBicubic{} of them, by \GapBicubicAbs{}
decibels on average; the gap is \NoiseMargin{} times the largest per-crop shift
the same model shows between two of its own runs on the same crops, so it is not
run-to-run variation. No perceptual measure was computed for those crops, which
is exactly why the comparison cannot be read as a verdict on the model: what it
establishes is that a fidelity-only readout ranks an interpolation above it, and
that reporting the model's \MethodPSNR{} decibels alone would have hidden that.
The third case is an earlier attempt over the full set of \CensusTotal{} inputs,
where a memory ceiling dropped \CensusSkipped{} of them, and the \CensusRan{}
that survived had no reference images and were never scored; which inputs
survived was fixed by output size rather than by design, so the measured set was
a property of the hardware. One recommendation follows from all three: print the
trivial arm beside the method, and treat any comparison the trivial arm wins as
a statement about the measurement rather than about the method.

\end{abstract}

%% The manuscript body, shared by every binding of this paper.
%% The class, the front matter and the bibliography style belong to the
%% binding; everything a reader reads belongs here, once.
\section{Introduction}

Every image restoration comparison has an arm that costs nothing to evaluate and
is almost never reported: the input, left alone. Bicubic upsampling is its
familiar cousin. Both are trivial in the sense that matters here --- no training,
no tuning, no choices --- and both are usually treated as too obvious to print.
The result is a table of methods with no floor under it.

This paper reports three comparisons from a single restoration project in which
the trivial arm was recorded anyway, and shows that in each of the three it
decided what the comparison could conclude. The comparisons differ in scale by
several orders of magnitude, in what was being compared, and in what went wrong.
What they share is that the trivial arm is the only reason any of that is
visible. Two of the three were then taken one step further, under rules written
down before the step was taken: the subset comparison was scored on a second,
perceptual endpoint, and the toy-scale stall was put through the initialisation
ablation the project had specified but never run. Both steps changed what the
comparison says, and in both cases the change is the result.

The first is a toy-scale run. A small convolutional restoration network and a
small transformer of comparable size were trained under a matched step budget and
compared with each other, and the project recorded that the transformer had lost.
The same record also holds the degraded input, scored as an arm: \ToyDegraded{}
decibels, against \ToyConv{} for the convolutional arm and \ToyTrans{} for the
transformer. Neither trained network reached the input it was asked to repair, so
there was no restoration to compare, and the ranking between the two was a
ranking of two failures. The project withdrew its own conclusion once this was
noticed, and the withdrawal is on disk with its reasons.

The second is the same convolutional arm at \BudgetMultiple{} times the budget.
The training objective fell by a factor of \ToyLossDrop{}. Fidelity against the
reference fell with it, by \ToyLongRegression{} decibels. Nothing about that pair
of numbers is visible from either one alone, and the trivial arm is what fixes their
interpretation: both budgets sit below the untouched input, so the curve is not a
learning curve approaching a target but a curve moving within a region where the
model is worse than doing nothing.

The third is a subset-scale evaluation of a released diffusion restoration
model~\cite{releasedmodel}. It ran to completion over \SubsetN{} crops with
reference images present, and the two interpolations scored on the same crops by
the same pass came out ahead: bicubic at \BicubicPSNR{} decibels and
nearest-neighbour at \NearestPSNR{}, against \MethodPSNR{} for the model. The
model is behind bicubic on \LossesBicubic{} of \SubsetN{} crops, by
\GapBicubicAbs{} decibels on average.

That third result is the one that has to be stated carefully, and the care is
most of the point. Peak signal-to-noise ratio (PSNR) measures distance from a reference.
A restoration model that synthesises plausible detail is expected to score worse
on it than one that hedges toward the conditional
mean~\cite{blau2018perception,ledig2017srgan}, and an interpolation is the
purest available hedge. So the fidelity ranking alone cannot say the model is
worse, and this paper's first draft said exactly that and stopped. The present
version does not stop there. Under a predeclared protocol, the same \SubsetN{}
crops were scored on a learned perceptual distance~\cite{zhang2018lpips}, and
the order reversed: the model is closer to the reference than bicubic on
\LpipsWins{} of \SubsetN{} crops, with \LpipsTies{} ties and an exact sign test
of \LpipsSignP{}, while losing to bicubic on PSNR on \PsnrLossesLocal{} of the
same crops. The two endpoints disagree on nearly every crop, in the direction
the perception--distortion trade-off predicts. A table carrying either endpoint
alone would have printed a ranking; the two together print the trade-off. Two
inference passes of the same model on the same crops differ by at most
\NoiseMax{} decibels per crop on PSNR. The record does not show the sampling
seed varied between the passes, so two runs cannot be read as an estimate of
sampling variance; the figure bounds only how far re-running the recorded
inference moved a crop. Reporting \MethodPSNR{} decibels without the
interpolation beside it would have presented that as a result; reporting the
interpolation's win without the perceptual endpoint beside it would have
presented the opposite.

A fourth thread runs underneath the third. An earlier census ran the same model
over the full set of \CensusTotal{} inputs and completed \CensusRan{} of them; a
memory budget skipped the other \CensusSkipped{} before inference, and the
\CensusRan{} that survived had no reference images and were never scored. Which
inputs survived that census was decided by output size, not by sampling. The
\SubsetN{} crops that were scored are the published RealSR crop128 protocol, not
the remnant of that skip rule. Weak-baseline problems and selection-by-infrastructure
problems are usually discussed separately; here they are two evaluations, and
only the census is a property of the hardware the run happened to have.

All of this is a finding about evaluation practice, with the architectures
serving as the material. The toy-scale stall, which the project had attributed
to zero-initialised residual gates, was put through the three-initialisation
ablation the project specified, and the ablation reversed the attribution: the
stall is initialisation-dependent, the arm that trains is the one whose gates
are exactly zero, and the two arms whose gates are not zero are the ones that
stall. The contribution is a demonstration, at three scales, that the cheapest
arm in the experiment is the one that determines whether the expensive arms
mean anything, together with the arithmetic showing exactly what it determined
in each case, a demonstration that the conclusion at one scale depends on which
endpoint is printed beside that arm, and a worked example of a recorded
attribution being retracted by the ablation written to test it.

Fig.~\ref{fig:scales} places the three comparisons on one page before they are
read separately. Every number on that page is restated later with its own
denominator; the overview exists so the shared role of the trivial arm is
visible at the outset.

\begin{figure}[tbp]
\centering
\includegraphics[width=\linewidth]{fig0_three_scales.pdf}
\caption{\CapScales}
\label{fig:scales}
\end{figure}

\section{Related work}

Two of this paper's results look, at first glance, like indictments of the
methods involved. Neither is, and the literature explains why in both cases. The
third result has less precedent, and is the one we think is most worth having.

The architectural setting is conventional. Super-resolution became a
deep-learning problem once a convolutional network could be trained end-to-end
against a
reference~\cite{dong2016srcnn,dong2016fsrcnn,kim2016vdsr,kim2016drcn,
shi2016espcn,wang2021survey}, then a residual one, then a residual-dense or
back-projection
one~\cite{he2016resnet,lim2017edsr,lai2017lapsrn,zhang2018rdn,zhang2017dncnn,
haris2018dbpn}. Adversarial and perceptual losses moved the same task toward
plausible texture~\cite{ledig2017srgan,johnson2016perceptual,wang2018esrgan},
and later work trained against synthetic degradations meant to stand in for
photographs~\cite{zhang2021bsrgan,wang2021realesrgan,agustsson2017div2k,
benchmarksubset}. Transformer restorers and latent diffusion restorers are the
current default
alternatives~\cite{vaswani2017attention,dosovitskiy2021vit,liu2021swin,
chen2021ipt,liang2021swinir,zamir2022restormer,ho2020ddpm,rombach2022ldm,
saharia2023sr3,saharia2022palette,releasedmodel}. The convolutional
encoder--decoder with skip connections remains the default dense-prediction
backbone~\cite{ronneberger2015unet}; batch, instance and layer normalisation,
then adaptive instance normalisation and feature conditioning layers, are the
usual way to inject a side
signal~\cite{ba2016layernorm,ulyanov2016instance,dumoulin2017style,
huang2017adain,perez2018film}; transformer backbones with adaptive layer
normalisation are now a standard alternative~\cite{peebles2023dit}. Our
toy-scale observation concerns the conditioning path of that last family and is
reported as a property of one small run, not of the architecture.

The metric result needs the most care. On the subset measured here an
interpolation is closer to the reference than the published method on
fidelity metrics, and read naively that is an embarrassing outcome for the
method. It is not, and the reason has been known since the perception--distortion
tradeoff was made precise: distortion and perceptual quality cannot both be
optimised without limit, so a method tuned for perceptual realism is
\emph{expected} to lose on squared-error-derived measures~\cite{blau2018perception}.
This was visible even before it was formalised, in the observation that
adversarially trained super-resolution produces more convincing images while
scoring worse on peak signal-to-noise
ratio~\cite{ledig2017srgan}. Structural similarity (SSIM)~\cite{wang2004ssim} is a
distortion measure in this sense and inherits the tradeoff; learned perceptual
distances~\cite{zhang2018lpips} sit on the other side of it. This paper
measures both sides on the same crops. On the distortion axis the interpolation
is closer to the reference on nearly every crop; on the perceptual axis the
released model is closer on nearly every crop. The honest reading of the subset
result is therefore that it locates the arms on both axes of a two-axis problem
and finds them in opposite orders, which is the trade-off observed rather than
a verdict on either arm.

Where we do claim a contribution is in evaluation hygiene, and there the
precedent is direct. Careful reproductions have repeatedly found that simple,
properly tuned baselines match or beat published neural methods once the
comparison is conducted on equal terms~\cite{dacrema2019worrying}, and the
broader argument that empirical progress is often an artifact of insufficiently
rigorous comparison is by now well
rehearsed~\cite{sculley2018winners,lin2019neural,ioannidis2005why,
simmons2011false,button2013power,cohen1992power,head2015phacking,
amrhein2019retire}. Our case adds a mechanism those critiques do not cover, and
it belongs to a separate earlier census rather than to the subset we scored.
The subset means reported here are the published RealSR crop128 protocol:
\SubsetN{} crops, with \EvalSkipped{} memory skips and \EvalReferences{} matched
references. The skip rule belongs to an earlier attempt over \CensusTotal{}
inputs. A memory ceiling determined which of those census inputs could be
processed, and the inputs it excluded were the larger ones from a particular
source. That census-level missingness is related to a property of the data
itself, which is precisely the regime in which ignoring the missing cases
biases the result~\cite{rubin1976inference}. A comparison conducted on the
survivors of a resource constraint is a comparison on a non-random
subsample---a selection effect on the census, not on the crop128
protocol---and the constraint is rarely reported because it feels like an
implementation detail rather than a design choice. We report the census of
what was excluded and the sizes behind each exclusion for that reason.

The statistical treatment is deliberately unambitious. Intervals come from the
bootstrap~\cite{efron1979bootstrap}, and paired comparisons use the sign
test~\cite{dixon1946sign}, the signed-rank test~\cite{wilcoxon1945individual}
and the paired $t$ test~\cite{student1908probable} in preference to any
aggregate that would collapse the per-crop structure the measurement produced.

\section{Methods}
\subsection{What was already on disk, and what was added under predeclaration}

The study has two layers. The first is a re-reading of records written by runs
that finished before this analysis began: three comparisons, none of which was
generated for this paper. The second is two computations made for this paper
under written predeclarations, each frozen and hashed before its first result
existed: a perceptual scoring pass over the subset crops, and a
three-initialisation ablation of the toy transformer. Every quantity in both
layers is tied to the manuscript by an evidence-manifest binding; the figure
code verifies each bound file before reading it, re-derives the summaries from
the per-crop or per-cell rows, holds each result to its own predeclaration, and
stops the build on any mismatch.

Three comparisons are read. The first is a toy-scale training run in which a
small convolutional restoration network and a small transformer of comparable
size were trained under a matched step budget on the same data, evaluated on the
same \ToyN{} held-out pairs at the same number of sampling steps, and scored
against the same degraded inputs. The second is the same convolutional arm
retrained at \BudgetMultiple{} times the budget, written outside the results tree
and never promoted. The third is a released diffusion restoration model run over
the \SubsetN{} crops of the published RealSR crop128 protocol, scored against
reference images, alongside bicubic and nearest-neighbour upsampling of the same
inputs.

\subsection{The trivial arms, and why they are arms}

Each comparison has an arm that performs no restoration. At toy scale it is the
degraded input passed through unchanged, which the training record scores
alongside the two trained networks. At subset scale there are two: bicubic and
nearest-neighbour upsampling of the same low-resolution inputs, scored by the
same metric pass over the same crops.

These are treated here as arms rather than as context. An arm that does nothing
is the weakest hypothesis a restoration comparison can be asked to beat, and a
method that does not beat it has not been shown to restore anything, whatever
its ranking against other methods. The three comparisons below differ in what
the trivial arm reveals, not in whether it should have been reported.

\subsection{What the fidelity measures are, and what they are not}

Two fidelity measures appear: peak signal-to-noise ratio, computed over the full
colour image and separately over the luma channel with a four-pixel border
removed, and structural similarity~\cite{wang2004ssim} under the same two
conventions. Both compare a restored image with a reference pixel by pixel or
window by window, and both reward outputs that stay close to the reference.

Neither is a measure of perceptual quality, and the difference matters for
everything reported at subset scale. Restoration methods that synthesise
plausible detail are known to score worse on distortion measures than methods
that produce a blurrier output closer to the conditional
mean~\cite{blau2018perception,ledig2017srgan}, and the trade-off is a property
of the estimator rather than a defect of any particular
model~\cite{blau2018perception}. A comparison with a distortion measure alone
can establish which output is closer to the reference; it cannot establish
which output is better, and a comparison with a perceptual measure alone cannot
either. This paper reports both, and treats their disagreement as the result.

\subsection{The perceptual endpoint}

The perceptual endpoint is the learned perceptual image patch similarity
(LPIPS), version 0.1, with the AlexNet backbone and its released linear
calibration head~\cite{zhang2018lpips}: lower is closer. Each restored output
and its reference are the full \(512\times512\) colour crop, scaled to
\([-1,1]\), with no border removed. The primary contrast for crop
\(i\in\mathcal{M}\) is \(\Delta^{\mathrm{LPIPS}}_i =
\mathrm{LPIPS}_{\mathrm{bicubic},i}-\mathrm{LPIPS}_{\mathrm{model},i}\), so a
positive value is a crop on which the model is perceptually closer. The
primary test is the exact two-sided sign test on the \SubsetN{} paired
contrasts at \(\alpha=0.05\), ties counted and excluded; a percentile bootstrap
of the mean contrast over 1,000 resamples and the per-camera win counts are
descriptive. The protocol, including these choices and the three possible
readings of the outcome, was written and hashed before any perceptual value was
computed, and the figure code re-derives the wins, the test and the mean from
the per-crop rows.

Two disclosures attach to this endpoint. First, provenance: the bound fidelity
table was produced by an earlier inference run on another machine, and its
restored pixels were not the ones scored here. The perceptual pass drew a fresh
sample of the released model on the same \SubsetN{} low-resolution inputs, at
sampling seed \LpipsSeed{} with \LpipsSteps{} sampling steps and the released
weights, and scored that sample. A predeclared anchor rule governs whether the
fresh sample may be read beside the bound ranking: its mean RGB PSNR must lie
within half a decibel of the bound record's, and it must lose to bicubic on at
least ninety of the \SubsetN{} crops; had either failed, the perceptual numbers
were to be reported as not comparable and no rewriting of the thesis was to
occur. The anchor passed with a shift of \AnchorDeltaDb{} decibels and
\AnchorLosses{} losses, and the bicubic arm recomputed for the pass reproduces
the bound bicubic rows to floating-point precision. The perceptual comparison
is therefore between the fresh sample and bicubic; the PSNR comparison in the
same figure uses the fresh sample as well, and the bound PSNR table is unchanged.

Second, the predeclaration was amended once. The predeclared scoring script
crashed while serialising its summary, after the per-crop values had been
computed and before any formal endpoint file was written. The amendment
replaced the script with a version that casts the offending values to plain
types --- the endpoints, pairing, test, threshold and anchor rule were not
changed --- and recorded the secondary backbone as not computed (below). The
primary result had by then been observed on a provisional side run that
imported the predeclared script's own functions; the amendment says so, and the
formal run under the amended protocol reproduced the provisional run exactly:
the same \LpipsWins{} wins, the same sign-test probability to the last digit,
the same anchor shift. Both the original predeclaration and the amendment are
bound and hashed with the result.

The protocol named LPIPS with the VGG16 backbone as a secondary robustness
check, never to be promoted. It was not computed: the cached VGG16 checkpoint
on the scoring machine was truncated, and the download route the workspace
permits was unavailable at the time of the run. Its status is recorded as such
in the bound summary, the per-crop rows carry null values for it, and nothing
is imputed for it. A distinct learned quality score that would have required a
further download was likewise omitted rather than substituted.

\subsection{The initialisation ablation}

The toy transformer was retrained under three initialisations with everything
else fixed: the library's default construction, which is the arm behind the
bound toy record; a standard transformer scheme applied to the same graph
without any zeroing step; and that scheme plus exact zeroing of the adaptive
normalisation gates and the final projection, which is the published
recipe~\cite{peebles2023dit}. Each was trained at \InitSeeds{} seeds for
\InitSteps{} steps on the same data, batch and learning rate as the bound run,
and evaluated on the same \ToyN{} held-out pairs, giving \InitCells{} cells. An
initialisation probe recorded, before training, whether each cell's gates were
exactly zero.

The predeclaration fixed the stall criterion before any cell ran: a cell
stalls if its mean objective over the last ten steps is at least
\InitThreshold{}, and a mode stalls if at least two of its three seeds do. Four
reading rules were written in evaluation order, the first match to fire: no
cell stalls; every cell stalls; the default stalls and at least one alternative
escapes; anything else. Restored fidelity against the degraded input is
reported per cell but does not enter the rule, and no cell beyond the nine was
to be added whatever the outcome. The frozen sources were re-hashed after the
run and match; the repository head advanced during the run because of
unrelated commits, which the receipt records, and the code hash stamped in
every cell is identical. The default cell at seed zero is the bound
construction at the bound seed, and the figure code stops if it fails to
reproduce the bound objective and fidelity within the printed precision.

The scored means and the memory census are two populations, not one union. Let
\(\mathcal{U}\) be the candidate input pool. The displayed estimator is the
published RealSR crop128 protocol with a matched reference:
\[
  \mathcal{M}
  = \{i\in\mathcal{U}: i\text{ is in the published crop128 protocol and
      has a matched reference}\},
\]
with \(N_{\mathcal{M}}=|\mathcal{M}|=\SubsetN{}\). That protocol run recorded
\(\EvalSkipped{}\) memory skips and \(\EvalReferences{}\) matched references. A
separate earlier census applied a memory filter over \(\CensusTotal{}\) inputs;
write
\[
  \mathcal{C}
  = \{i\in\mathcal{U}: i\text{ was presented to that memory-filtered
      census}\}.
\]
The displayed estimator does not take \(\mathcal{M}\cup\mathcal{C}\) or
intersect \(\mathcal{M}\) with the memory filter. All subset means, paired
differences, intervals, and win counts are functions of \(\mathcal{M}\) only.
An input skipped for memory in the census is counted in \(\mathcal{C}\) and is
outside \(\mathcal{M}\); it contributes a census count, not a zero-valued
metric. Thus the reported comparison is the crop128 protocol and cannot be
read as an estimate for \(\mathcal{C}\) or for the released model on a full
benchmark.

\subsection{Checks the build enforces}

Several sentences in this paper are enforced as build-time checks rather than
asserted in prose, and the figure code stops rather than reporting a number that
would contradict them.

The one-line ladder entries that summarise the toy run must still restate the
values the training record holds. The record that withdrew the toy comparison
must agree with the training record on both parameter counts and with the
unpromoted longer run on its restored fidelity. Every mean printed in the subset
summary must be the mean of the per-crop rows it claims to summarise, and every
recorded interval must be centred on that same mean. The recorded paired
difference, its extremes, and the recorded win and loss counts, overall and
within each source group, must all reproduce from the per-crop rows. The per-crop
rows themselves are joined across arms on the crop name rather than on position,
so a reordering cannot silently pair the wrong crops.

Two recorded paths hold the record for the bicubic arm. Comparing their manifest bindings
establishes whether that is a second measurement or a second copy, and the
manuscript reports which; comparing their rows establishes that they still agree
either way.

\subsection{Bounding the model against itself}

Two pixel sets exist for the released model over the same \SubsetN{} crops,
produced by two inference passes and scored by the same protocol. The per-crop
differences between them are not used as a result. They are used once, to bound
how far re-running the recorded inference moved a crop. The record does not show
the sampling seed varied between the passes, so two runs cannot be read as an
estimate of sampling variance. The largest of the differences is drawn on the
figure that reports the gap, and it is reported as a largest observed shift
rather than as an interval.

\subsection{Tests}

The paired difference between the released model and bicubic upsampling is
tested three ways on the same \SubsetN{} crops: the exact two-sided sign
test~\cite{dixon1946sign}, which counts how many crops fall each way and assumes
nothing about their distribution; the Wilcoxon signed-rank
test~\cite{wilcoxon1945individual} on the same pairs; and the paired $t$
test~\cite{student1908probable} with its interval, which is the only one of the
three that assumes a distributional shape. A percentile
bootstrap~\cite{efron1979bootstrap} over the paired differences is reported
beside the interval each record already carries, so the two constructions can be
compared rather than one standing in for the other. Source group is a recorded
property of each crop and is reported as a stratum, not as a covariate: the
strata were not chosen after seeing the differences, and no adjustment is made.
The resampling unit is the crop, so these intervals describe this measured
subset and do not provide camera-level or population-level uncertainty when
crops from the same source are correlated.

For a crop \(i\in\mathcal{M}\), the fidelity contrast is
\(\Delta_i=\mathrm{PSNR}_{\mathrm{model},i}
-\mathrm{PSNR}_{\mathrm{bicubic},i}\) (and analogously for SSIM). These are
distortion contrasts, not latent perceptual-quality contrasts; a negative
\(\overline{\Delta}\) says only that the model is farther from the available
reference under that metric.

The cross-metric consistency panel computes this contrast separately for each
of RGB PSNR, luma PSNR, RGB SSIM and luma SSIM before taking the paired mean.
Each interval is the 2.5th and 97.5th percentiles of 1,000 resamples of the
corresponding paired differences. PSNR remains in decibels and SSIM remains
unitless, so the two metric families are displayed on separate native axes
rather than combined into a standardized score. This panel is descriptive: it
does not add a multiplicity-adjusted test or turn distortion differences into a
claim about perceptual quality.

\subsection{The census of an attempt, and the path it names}

An earlier attempt ran the same released model over the full set of
\CensusTotal{} inputs under a fixed memory budget, expressed as a maximum number
of output pixels and a maximum output side length. Inputs whose upsampled output
would exceed the budget were skipped before inference. The counts, the budget,
the size distribution of what was skipped and the source-group composition of
what was skipped are all recorded.

The census names a live path as its source, and that path \CensusPathState{}:
the later subset run overwrote it. A snapshot of the earlier state was retained
under a history directory, and the census is checked against that snapshot
rather than against the live path. The check is what makes the census usable;
without the retained snapshot the counts would be unverifiable, and the paper
reports the near miss rather than passing over it.

\subsection{The boundary of what is computed}

The perceptual computation is the single endpoint described above, and the
secondary backbone the protocol named is reported as not computed rather than
estimated. The \CensusRan{} outputs of the earlier attempt enter only as a
count: they have no reference images on disk, so they carry no value, and the
build stops if any of them ever acquires one, because their absence is the
reason those outputs cannot be scored. The two larger comparisons this line of
work planned stay at their recorded status --- \PrimaryTier{} for the
full-length reference comparison and \SotaTier{} for the comparison against the
strongest published method that would follow it --- and the subset is reported
as the subset, with completing it leaving both where they were. The scored
\SubsetN{} crops are the published crop128 protocol; the memory skip belongs to
the earlier census.


\section{At toy scale, no arm reached the input}

Table~\ref{tab:toy} and Fig.~\ref{fig:toy} give the toy-scale run as its own
records describe it. Two networks were trained for \ToySteps{} steps on the same
data at seed \ToySeed{}, and both were evaluated on the same \ToyN{} held-out
pairs at \ToyEvalSteps{} sampling steps against the same degraded inputs. The
convolutional arm~\cite{ronneberger2015unet} reached \ToyConv{} decibels. The
transformer arm~\cite{peebles2023dit} reached \ToyTrans{}. The degraded input,
scored by the same pass, is at \ToyDegraded{}.

The ordering of the two trained arms is the comparison the project originally
recorded. The position of both relative to the third number is what makes that
comparison unavailable. A restoration method that scores below its own input has
not restored; it has degraded further. Ranking two such methods against each
other ranks the sizes of two failures, and the shortfall in question is not
marginal: \ToyConvShortfall{} decibels for the better arm and
\ToyTransShortfall{} for the worse. The structural similarity values say the same
thing from the other direction, \ToyConvSSIM{} and \ToyTransSSIM{} against
\ToyDegradedSSIM{} for the untouched input.

The project reached the same conclusion independently and recorded it. Its
withdrawal note states that the earlier wording is not supportable and must not
be quoted, and lists both the reason above and several others. We take that
record as the starting point rather than as a result of ours, and the figure code
checks that the withdrawal still agrees with the training record it describes.

Two of the other reasons are worth reporting because they generalise. The first
is that the shared budget was a step budget and nothing else. The two arms differ
by \ToyParamExcess{} per cent in trainable parameters, \ToyTransParams{} against
\ToyConvParams{}, so they were not matched on capacity. They also differ by a
factor of \ToySpeedRatio{} in wall clock for the same \ToySteps{} steps,
\ToyTransSeconds{} seconds against \ToyConvSeconds{}, so they were not matched on
compute either. A budget matched on steps, on parameters and on wall clock are
three different experiments, and at this ratio the first two would give opposite
answers to the third. Naming which one is being held fixed is not pedantry when
the gap between them is thirtyfold.

The second is that the transformer's number is not a measurement of the
transformer. Its training objective was \ToyTransLoss{} at the last step, with a
mean of \ToyTransLastTen{} over the last ten --- flat, at unity, having not
moved. A denoising objective normalised against unit-variance noise sits at
approximately one when the network contributes nothing~\cite{ho2020ddpm}, so
\ToyTrans{} decibels measures a network that never left its initialisation, and
the number carries no information about the architecture at any budget. The
project's record had attributed the stall to the zero-initialised residual
gates that the transformer's adaptive normalisation layers use at
initialisation~\cite{peebles2023dit}. That attribution was testable, and the
project's own design said how: a three-way comparison across the library's
default initialisation, a standard scheme without the zeroing step, and the
scheme that zeroes the gates and final projection exactly, with the training
objective and restored fidelity as endpoints.

That ablation was run for this paper, under a predeclaration that fixed the
stall criterion --- a mean objective of at least \InitThreshold{} over the last
ten steps --- and the reading rules before any cell was trained
(Fig.~\ref{fig:init-ablation}). Nine cells were trained: the three
initialisations at \InitSeeds{} seeds each, for the same \InitSteps{} steps on
the same data. The library default stalls in \InitDefaultStall{} of
\InitSeeds{} seeds, with a mean last-ten objective of \InitDefaultLoss{}; the
standard scheme stalls in \InitStandardStall{} of \InitSeeds{}, at
\InitStandardLoss{}. The zero-gated scheme escapes in \InitZeroEscape{} of
\InitSeeds{}: its objective falls to \InitZeroLoss{}, and its restored fidelity
is \InitZeroPsnr{} decibels against \InitStallPsnr{} for the two arms that
stalled. The seed-zero default cell reproduces the bound record --- the same
construction, seed and budget --- to within the printed precision, so the
stall is a property of that construction and not of one run.

So the stall is initialisation-dependent, and the project's attribution of it
is contradicted rather than confirmed. An initialisation probe recorded before
training shows why: the library default does not zero the gates, the standard
scheme does not zero them, and the one arm that does zero them exactly is the
arm that trains. Whatever produced the stall in the default construction, it
was not the zero-initialised gates, because zeroing the gates is what removes
it. We retract the attribution as it stood in the earlier draft and report the
ablation in its place. What does not change is the trivial arm's verdict: every
one of the nine cells, including the three that trained, finishes below the
degraded input at \IdentityPsnr{} decibels, so the ablation moves the toy
transformer from a network that did not train to a network that trained and
still did not restore, and nothing about the architecture follows from either.

\begin{figure}[tbp]
\centering
\includegraphics[width=\linewidth]{fig11_init_ablation.pdf}
\caption{\CapInitAblation}
\label{fig:init-ablation}
\end{figure}

\begin{table}[tbp]
\centering
\caption{\CapToyTable}
\label{tab:toy}
\small
% Generated by the figure script from hash-bound evidence. Do not hand-edit.
\begin{tabular}{lrrrrr}
\toprule
Arm & Parameters & Steps & Wall clock (s) & Final loss & PSNR (dB) \\
\midrule
Degraded input, left alone & --- & --- & --- & --- & 20.75 \\
Convolutional arm, 50 steps & 16,058,179 & 50 & 47.46 & 0.0810 & 14.72 \\
Convolutional arm, 500 steps & 16,058,179 & 500 & 478.17 & 0.0338 & 14.05 \\
Transformer arm, 50 steps & 17,961,024 & 50 & 1.58 & 0.9912 & 6.57 \\
\bottomrule
\end{tabular}

\end{table}

\begin{figure}[tbp]
\centering
\includegraphics[width=\linewidth]{fig1_toy.pdf}
\caption{\CapToy}
\label{fig:toy}
\end{figure}

\section{The longer run, read on both axes}

Fig.~\ref{fig:budget} adds the one arm that tests whether the toy shortfall was
a budget problem. The convolutional network was retrained at \ToyLongSteps{}
steps, \BudgetMultiple{} times the original, on the same data at the same
learning rate and seed, and evaluated identically.

The training objective did what more training is supposed to do: \ToyConvLoss{}
at \ToySteps{} steps, \ToyLongLoss{} at \ToyLongSteps{}, a fall by a factor of
\ToyLossDrop{}. Fidelity against the reference went the other way, from
\ToyConv{} decibels to \ToyLong{}, a loss of \ToyLongRegression{} decibels for
\ToyLongSeconds{} seconds of additional wall clock. Structural similarity rose
slightly over the same interval, from \ToyConvSSIM{} to \ToyLongSSIM{}, so the
two fidelity measures disagree about the direction as well.

Each of the three panels of that observation is unremarkable alone. A falling
objective is expected. A metric that moves against the objective is a familiar
kind of overfitting or objective mismatch. Two fidelity measures disagreeing is
ordinary. What makes the combination interpretable is the fourth number: both
budgets are below the untouched input, by \ToyConvShortfall{} decibels and more.
The model is not converging toward a target and overshooting. It is moving inside
a region in which its output is worse than its input, and the direction of travel
inside that region is not informative about what a larger budget would do.

This is the cheapest possible test of a budget explanation, and it is negative.
It also explains the record's location: the longer run was written outside the
results tree and was never promoted into the comparison ladder. It contradicts
the promoted run, and it is recorded here for exactly that reason. A record that
disagrees with the one that was kept is the record most worth keeping.

\begin{figure}[tbp]
\centering
\includegraphics[width=\linewidth]{fig2_budget.pdf}
\caption{\CapBudget}
\label{fig:budget}
\end{figure}

\section{At subset scale, an interpolation is closer on fidelity and the model is closer on perception}

The subset evaluation is a different kind of object. A released diffusion
restoration model built on a latent text-to-image
backbone~\cite{rombach2022ldm,releasedmodel} was run over \SubsetN{} crops from a
real-world paired super-resolution benchmark~\cite{benchmarksubset}, drawn in
equal numbers from \SubsetCameras{} camera groups, with reference images present
for all of them. The run completed: \EvalRun{} inputs processed, \EvalSkipped{}
skipped, \EvalReferences{} references matched. A separate metric pass scored the
model's outputs and, on the same crops, bicubic and nearest-neighbour upsampling
of the same low-resolution inputs.

Table~\ref{tab:subset} gives all four recorded measures. The model is at
\MethodPSNR{} decibels with a recorded interval from \MethodLow{} to
\MethodHigh{}. Bicubic upsampling is at \BicubicPSNR{} and nearest-neighbour at
\NearestPSNR{}. The ordering is the same on the luma channel with a border
removed, \MethodYPSNR{} against \BicubicYPSNR{}, and the same on structural
similarity, \MethodSSIM{} against \BicubicSSIM{} for bicubic and
\NearestSSIM{} for nearest-neighbour, a paired difference against the first of
those of \GapSSIM{}. All four measures agree, which
is not surprising --- they are four views of the same kind of quantity --- but it
does remove the possibility that the ordering is an artifact of one convention.

Fig.~\ref{fig:subset} plots the comparison crop by crop. The model is behind
bicubic on \LossesBicubic{} of \SubsetN{} crops and behind nearest-neighbour on
\LossesNearest{}. The mean paired difference against bicubic is \GapBicubic{}
decibels, with a percentile bootstrap interval from \GapLower{} to \GapUpper{}
recomputed here and a recorded interval from \GapRecordedLower{} to
\GapRecordedUpper{}, which agree to within the resampling; against
nearest-neighbour it is \GapNearest{}. The paired $t$ interval
is \GapTLower{} to \GapTUpper{}. The exact sign test returns \GapSignP{} and
Wilcoxon \GapWilcoxP{}; against nearest-neighbour the sign test returns
\GapNearestSignP{}. The spread of per-crop differences is wide, from
\GapWorst{} to \GapBest{} decibels, so the ordering is not uniform even though it
is nearly unanimous.

The fidelity ranking on its own does not say the model is worse, and the reason
is not caution for its own sake. Distortion measures and perceptual quality
trade off against each other in a way that is a property of the estimation
problem rather than of any particular model~\cite{blau2018perception}. A method
that samples plausible high-frequency detail will place that detail somewhere,
and where it places it will usually not be where the reference has it; the
resulting penalty is real and is not evidence of a defect. This has been
understood in super-resolution since adversarial methods first traded
distortion for realism~\cite{ledig2017srgan}, and it is why perceptual measures
were introduced~\cite{zhang2018lpips}. An earlier draft of this paper stopped
at that point, because the record then held no perceptual measure for these
crops. The measure has since been computed, under a protocol written down and
frozen before any value was seen.

The perceptual endpoint is the learned perceptual image patch similarity with
the AlexNet backbone~\cite{zhang2018lpips}, scored on the full crop against the
same reference, lower being closer. It was scored on a fresh local sample of
the released model, drawn at sampling seed \LpipsSeed{}, rather than on the
pixel set behind Table~\ref{tab:subset}, which an earlier run had produced with
the vendor's default seed; the predeclared anchor rule under which the fresh
sample may be read beside the bound ranking, and the disclosure that attaches
to the protocol, are in Methods. The anchor passed: the fresh sample's
mean PSNR is \AnchorLocalPsnr{} decibels, \AnchorDeltaDb{} above the bound
\MethodPSNR{}, and it loses to bicubic on \PsnrLossesLocal{} of \SubsetN{}
crops against the bound record's \LossesBicubic{}.

On that endpoint the order reverses (Fig.~\ref{fig:perceptual-reversal}). The
model's mean perceptual distance is \LpipsDit{} against \LpipsBicubic{} for
bicubic; the model is closer on \LpipsWins{} of \SubsetN{} crops and bicubic on
\LpipsLosses{}, with \LpipsTies{} ties, and the exact sign test returns
\LpipsSignP{}. The mean paired difference in bicubic's disfavour is
\LpipsDelta{}, with a percentile bootstrap interval from \LpipsDeltaLo{} to
\LpipsDeltaHi{}. The reversal holds within each camera group, \LpipsCanonWins{}
and \LpipsNikonWins{} wins of fifty. Read crop by crop, the two endpoints
disagree almost everywhere: on \QuadLossWin{} crops bicubic is closer on PSNR
and the model is closer on the perceptual distance, on \QuadLossLoss{} bicubic
is closer on both, on \QuadWinWin{} the model is closer on both, and on
\QuadWinLoss{} the model is closer on PSNR while bicubic is closer perceptually.
The secondary backbone the protocol named, VGG16, was later scored on the
same crops as a backbone-sensitivity check: the model is closer on
\LpipsVggWins{} of \SubsetN{} and bicubic on \LpipsVggLosses{}, sign test
\LpipsVggSignP{}, which does not meet the predeclared threshold. That check
does not confirm the AlexNet ranking and does not rewrite it; it is not
promoted, and the thesis remains the reversal on the AlexNet endpoint.

What the subset therefore establishes is a statement about measurement, not a
verdict: scored on fidelity, this evaluation ranks two interpolations above the
model on nearly every crop; scored on the perceptual endpoint, it ranks the
model above bicubic on nearly every crop; and the crops on which the two
endpoints agree are the minority. Two consequences follow. The first is that
neither endpoint alone adjudicates the model. The fidelity table would have
been read as an indictment and the perceptual table as a vindication, and both
readings would be readings of the measure. The second is practical. A table
containing \MethodPSNR{} decibels and no interpolation would be read as a
result; a table containing \BicubicPSNR{} beside it and no perceptual endpoint
would be read as the opposite result; each costs one line of code, and the
table means something different after each line is added.

\begin{table}[tbp]
\centering
\caption{\CapSubsetTable}
\label{tab:subset}
\small
% Generated by the figure script from hash-bound evidence. Do not hand-edit.
\begin{tabular}{lrrrr}
\toprule
Arm & PSNR (dB) & SSIM & Luma PSNR (dB) & Luma SSIM \\
\midrule
Nearest-neighbour upsampling & 24.65 & 0.674 & 26.23 & 0.712 \\
Bicubic upsampling & 25.48 & 0.709 & 27.11 & 0.746 \\
Published restoration method & 21.88 & 0.614 & 23.44 & 0.658 \\
\midrule
Method minus bicubic, paired & $-3.60$ & $-0.095$ & $-3.68$ & $-0.088$ \\
Method minus nearest, paired & $-2.77$ & $-0.060$ & $-2.79$ & $-0.054$ \\
\bottomrule
\end{tabular}

\end{table}

\begin{figure}[tbp]
\centering
\includegraphics[width=0.62\linewidth]{fig3_subset.pdf}
\caption{\CapSubset}
\label{fig:subset}
\end{figure}

\begin{figure}[tbp]
\centering
\includegraphics[width=\linewidth]{fig10_perceptual_reversal.pdf}
\caption{\CapPerceptualReversal}
\label{fig:perceptual-reversal}
\end{figure}

\section{What the gap is, and what it is not}

A gap of \GapBicubicAbs{} decibels invites two easy dismissals. Both can be
taken to the records rather than argued about; the records answer the second
fully and the first only in part.

The first is that it might be run-to-run variation. Two pixel sets exist for the
model over the same \SubsetN{} crops: two inference passes of the same model,
scored by the same protocol. The largest per-crop difference between them is
\NoiseMax{} decibels, the mean absolute difference is \NoiseMean{}, and the shift
in the mean is \NoiseShift{}. The record does not show the sampling seed varied
between the passes, so two runs cannot be read as an estimate of sampling
variance; the figure bounds only how far re-running the recorded inference moved
a crop, and it is reported as a largest observed shift rather than an interval.
What it rules out is correspondingly narrow. The gap is not an artefact of
re-running the recorded inference; how far a different seed would move the
model's output is not measured here. Fig.~\ref{fig:paired} draws that shift as a
band so the two scales can be seen against each other rather than quoted.

The second is that the gap might be one camera group, one crop, or one tail. It
is not. Within each of the \SubsetCameras{} groups separately the mean difference
is negative with a bootstrap interval excluding zero, and the group means differ
from each other by less than either differs from zero. The groups are a recorded
property of each crop, fixed before the differences were computed, so this is a
stratification and not a search. The model does win on \WinsBicubic{} crops
against bicubic and \WinsNearest{} against nearest-neighbour, and the best of
those is \GapBest{} decibels, so the ordering is not a floor effect either.

One provenance detail belongs here rather than in a footnote. Two recorded paths
hold the bicubic arm, and they are \BicubicRecordState{}. The second path is
therefore not an independent confirmation of the first, and this paper does not
treat it as one; it is reported because a reader counting records would otherwise
count two. The rows were compared as well as the manifest bindings, so the build would have
stopped had they diverged.

\begin{figure}[tbp]
\centering
\includegraphics[width=\linewidth]{fig4_paired.pdf}
\caption{\CapPaired}
\label{fig:paired}
\end{figure}

The same paired join can be checked without changing the population or quietly
putting unlike units on one axis. Fig.~\ref{fig:metric-consistency} repeats the
DiT--bicubic contrast for RGB and luma PSNR and for RGB and luma SSIM. PSNR is
shown in decibels and SSIM on its unitless scale in separate facets; each point
is the paired mean and each whisker is a percentile bootstrap interval over the
same \SubsetN{} crop128 pairs. This is a descriptive consistency check across
distortion measures, not a perceptual endpoint or an inference for the skipped
inputs.

\begin{figure}[tbp]
\centering
\includegraphics[width=\linewidth]{fig6_metric_consistency.pdf}
\caption{\CapMetricConsistency}
\label{fig:metric-consistency}
\end{figure}

The camera split is a sensitivity analysis, not a second population. Fig.~\ref{fig:camera-sensitivity}
shows the same paired RGB-PSNR contrast within each recorded camera stratum,
with the win count printed beside the interval. Both strata point in the same
direction, but the intervals describe the measured crops only: scene-level
clusters and the skipped census are not recoverable from these rows.

\begin{figure}[tbp]
\centering
\includegraphics[width=\linewidth]{fig7_camera_sensitivity.pdf}
\caption{\CapCameraSensitivity}
\label{fig:camera-sensitivity}
\end{figure}

Fig.~\ref{fig:crop-camera-metrics} keeps the crop-level PSNR and SSIM rows,
their Canon/Nikon descriptive intervals, and the earlier \CensusTotal{}-input
census in one scope-aware diagnostic. Its intervals are recomputed from the
bound crop rows; the bound camera-stratum record supplies the camera counts and means. It
does not add a perceptual metric or promote the measured subset to the official
full table.

\begin{figure}[tbp]
\centering
\includegraphics[width=\linewidth]{fig9_crop_camera_metrics.pdf}
\caption{\CapCropCameraMetrics}
\label{fig:crop-camera-metrics}
\end{figure}

\section{A memory ceiling chose an earlier census, not the scored crops}

The \SubsetN{} crops are the published RealSR crop128 protocol. They are not
what remained after a constraint that had nothing to do with the science. An
earlier census of a mixed input pool is what the memory rule selected, and the
record of how that happened is complete enough to reconstruct.

An earlier attempt ran the same model over the full set of \CensusTotal{} inputs
under a fixed memory budget: at most \BudgetPixels{} output pixels and at most
\BudgetLong{} pixels on the long side, chosen for the \BudgetVram{}-gibibyte
accelerator the run had. Inputs whose upsampled output would exceed the budget
were skipped before inference rather than attempted and recovered. Of
\CensusTotal{} inputs, \CensusSkipped{} were skipped, or \CensusSkipShare{} per
cent, and \CensusRan{} ran.

The skip rule was deterministic and single-valued. Every skipped input was of one
size class, whose upsampled output would have been \SkippedSide{} pixels on a
side, \SkippedPixels{} pixels in total, \BudgetExcess{} times the budget. Nothing
was skipped for being hard, or unusual, or badly exposed. The evaluation set was
partitioned by output dimensions, and one side of the partition was discarded.

Table~\ref{tab:census} and Fig.~\ref{fig:ceiling} give what that removed. The
dropped inputs span \CensusCameras{} source groups, of which \AbsentGroups{} do
not appear anywhere in the measured subset and so are absent from every number in
this study. Whether this matters depends on whether restoration difficulty varies
by capture device, which these records cannot answer; the honest statement is
that the question is unanswerable from this evaluation and that the evaluation
does not look like it was designed that way.

The \CensusRan{} inputs that did survive make the point sharper. They were
processed, their outputs were written, and none of them was ever scored, because
no reference image exists on disk for any of them. The build stops if one ever
acquires one. So the earlier attempt produced \OrphanStems{} outputs and zero
measurements: the ceiling did not merely shrink the evaluation set, it left the
survivors unmeasurable, and the two failures compounded.

What happened next is the reasonable response and is worth recording as such.
The scored comparison was taken on the published crop128 protocol, whose
\SubsetN{} crops fit the budget and come with reference images. That protocol
is not the remnant of the census skip rule. The decision is sound and the
resulting run is complete. The full comparison on the
released split was never produced (its recorded status is \PrimaryTier{}), nor was
the comparison against the strongest published method that would follow it
(\SotaTier{}), and completing the subset changes neither.

One near miss is worth reporting because it very nearly removed the ability to
say any of this. The census that counts \CensusSkipped{} skipped inputs cites a
live path as its source, and that path \CensusPathState{}: the later subset run
overwrote it, and the file now at that location records \EvalRun{} inputs run and
\EvalSkipped{} skipped. A snapshot of the earlier state was retained under a
history directory, and the census reproduces exactly against the snapshot. The
mitigation worked. The mechanism that made it necessary --- a derived table
pointing at a mutable path rather than at an immutable one --- is one line of
schema, and it is the difference between a checkable count and an anecdote.

\begin{table}[tbp]
\centering
\caption{\CapCensusTable}
\label{tab:census}
\small
% Generated by the figure script from hash-bound evidence. Do not hand-edit.
\begin{tabular}{lrl}
\toprule
Source group & Inputs dropped & Present in the measured subset \\
\midrule
Canon & 55 & yes \\
Nikon & 50 & yes \\
DSC & 20 & no \\
panasonic & 20 & no \\
sony & 17 & no \\
IMG & 12 & no \\
P114 & 12 & no \\
P116 & 4 & no \\
P117 & 3 & no \\
\midrule
All groups & 193 & --- \\
\bottomrule
\end{tabular}

\end{table}

\begin{figure}[tbp]
\centering
\includegraphics[width=\linewidth]{fig5_ceiling.pdf}
\caption{\CapCeiling}
\label{fig:ceiling}
\end{figure}

\section{Discussion}

The three comparisons are not variations on one mistake. They fail in different
ways, and the trivial arm plays a different role in each: at toy scale it
invalidates the comparison, at extended budget it fixes the interpretation of a
trend, and at subset scale it reveals that the measure in use decides the
ranking, because the second measure reverses it. What travels is the practice,
not the finding.

A method below the trivial arm has no rank. The strongest form of the toy-scale
result is also the simplest: two methods below the input they were given are
not ordered by their distance below it. This is an admissibility condition
rather than a statistical subtlety, it needs no test, and it costs one number
to check. Comparisons against weak or absent baselines have been shown to
overturn reported progress in several
fields~\cite{dacrema2019worrying,lin2019neural,sculley2018winners}; the case here
is the limiting one, where the baseline is the arm that does nothing, and it
still wins.

A trend needs a level. Raising the budget \BudgetFold{} improved the objective
and worsened the output. Either half of that, alone, has a standard reading,
and the level is what settles which reading applies: the whole trajectory lies
in a region where the model is worse than its input. Learning curves are
reported constantly and are almost never reported with a reference level on the
same axis. Adding one is free.

The trivial arm is a diagnostic for the measure as much as for the method. At
subset scale the interpolation's fidelity win tells us that the measure in use
ranks a hedge above a sampler, which is what that measure does. The earlier
draft of this paper stopped at that signal and stated the rule it implies: when
the trivial arm wins by a wide margin on a distortion measure, the next action
is to compute a perceptual measure before publishing the ranking. This version
followed the rule, under a protocol frozen first, and the perceptual measure
reversed the ranking on \LpipsWins{} of \SubsetN{} crops. The rule therefore
has a demonstrated consequence as well as a rationale: had the fidelity table
been published alone, its ranking would have been inverted by the next
endpoint. The trivial arm is still what made the reversal legible, because a
perceptual score for the model with no interpolation beside it would have been
a number rather than an ordering.

An attribution is a hypothesis until the ablation runs. The toy-scale stall had
a plausible mechanism attached to it, and the ablation found the mechanism
pointed the wrong way. Testing it cost nine short runs. The general point is
about records rather than gates: a recorded explanation of a failed run tends
to outlive the run, and the cheapest way to keep it from becoming a finding is
to specify the ablation when the explanation is written and to run it before
the explanation is printed.

Infrastructure limits are selection effects. A memory budget that skips large
inputs is a non-random selection rule applied to the census population, and the
resulting sample is informative about the part of that pool that fits. The
scored crop128 comparison is a different population: a published protocol, not
the remnant of the skip. Applied to the census, this has the same structure as
any missing-not-at-random problem~\cite{rubin1976inference}, with the
difference that the mechanism is fully known and recorded, which is a
considerable advantage rarely available elsewhere. Reporting the skip rule, its
threshold, the counts and the composition of what was dropped takes a
paragraph, and it converts an unstated restriction of scope into a stated one.

Withdrawal should be as citable as publication. The most useful single artifact
in this study is the record in which the project withdrew its own toy
comparison, listed six reasons, and wrote down the narrower statement its data
still supported. Nothing was deleted. The original claim, the withdrawal, the
reasons and the replacement all sit in one file, and the contradicting longer
run sits beside it. That is what made a second reading possible at all, and it
is a practice any project can adopt without a tool.

Derived tables should cite immutable sources. The census survived only because
a snapshot was retained. Derived summaries that name a live path are correct
when written and become unverifiable the moment the path is rewritten, in a way
that leaves no trace in the summary itself. Naming an evidence ID instead of a
path is a schema change rather than a process change.

\section{Scope of the results}

The subset result is a statement about two endpoints on one set of crops. On
the published RealSR crop128 protocol the interpolations are closer to the
reference on the distortion measures and the released model is closer on the
learned perceptual distance that was computed; the two endpoints order the arms
oppositely, and this paper reports that disagreement as the finding. It holds
for one learned perceptual distance with one backbone, scored on these
\SubsetN{} crops with no reader study beside it; the secondary backbone the
protocol named was later scored as a backbone-sensitivity check and is not
promoted, and the result describes what this
evaluation establishes under each endpoint rather than the model's quality.

The crops measured are the published crop128 protocol, and the two larger
comparisons the project planned remain at their recorded status: the full
comparison on the released split is \PrimaryTier{} and the comparison against
the strongest published method that would follow it is \SotaTier{}. Completing
the subset leaves both where they were, and every number in this paper is a
number about the crops that were scored or the census that was counted.

The toy-scale results describe one small transformer at one budget. Under the
library's default initialisation its objective stayed at unity in the bound run
and in all \InitSeeds{} ablation seeds; under the initialisation that trained,
it finished below the degraded input. A toy run that did not train, and a toy
run that trained and did not restore, locate this construction at this budget
and say nothing about diffusion transformers as an architecture at any other
scale; nothing here connects the toy transformer to the released model.

The attribution of the stall to zero-initialised residual gates, relayed in the
earlier draft of this paper from the project's record, is withdrawn: the
predeclared ablation found that the arm whose gates are exactly zero is the arm
that trains, and the two arms whose gates are not zero are the arms that stall.
What the ablation supports is the narrower statement that on this toy, at this
budget, the stall depends on the initialisation and zeroing the gates removes
it. The mechanism behind the default construction's stall is left open, and the
published initialisation scheme is characterised only on this toy.

The perceptual computation that carries the thesis is the one named in
Methods, on the \SubsetN{} crops. The VGG16 backbone the protocol named was
later scored as a backbone-sensitivity secondary and is not promoted; no
distortion-free quality score stands in for the AlexNet endpoint, and every
perceptual value in this paper belongs to one of those crops.

The \CensusRan{} outputs of the earlier attempt appear only as a count. They
have no reference images on disk, so they carry no measurement, and the figure
code stops if any of them ever acquires a reference, because their absence is
the reason they cannot be scored. The \CensusSkipped{} skipped inputs were
never processed, so what is known about them is their size and their source
group, and the statement made about them is about the selection rule that
removed them.

\section{Conclusion}

Three comparisons from one restoration project were re-read against the arm that
does nothing, and two of them were then taken one predeclared step further. At
toy scale, two trained networks finished at \ToyConv{} and \ToyTrans{} decibels
against \ToyDegraded{} for the untouched input, so the comparison between them
ranked two failures and the project withdrew it. The transformer's stall
reproduced in \InitDefaultStall{} of \InitSeeds{} seeds under the default
initialisation and vanished in \InitZeroEscape{} of \InitSeeds{} under the
scheme that zeroes the gates exactly, so the stall is initialisation-dependent
and the recorded attribution of it to zero-initialised gates is withdrawn; every
cell still finished below the input. Raising the training budget \BudgetFold{}
on the better arm cut the objective by a factor of \ToyLossDrop{} and moved
fidelity \ToyLongRegression{} decibels the wrong way, with both budgets below
the untouched input. At subset scale, a released restoration model scored
\MethodPSNR{} decibels against \BicubicPSNR{} for bicubic upsampling on the same
\SubsetN{} crops, losing on \LossesBicubic{} of them by \GapBicubicAbs{}
decibels on average; on a learned perceptual distance scored on the same crops
it is closer to the reference on \LpipsWins{} of them, with a sign test of
\LpipsSignP{}. Re-running the recorded inference moved no crop by more than
\NoiseMax{} decibels on PSNR; because the record does not show the sampling
seed varied between the two passes, that number bounds only what a re-run moved
and is not an estimate of sampling variance.

That last comparison is not a verdict on the model in either direction. It is a
verdict on the evaluation: measured on fidelity it ranks an interpolation first,
measured on the perceptual endpoint it ranks the model first, and a table
reporting either ranking alone would have concealed the other. The scored
\SubsetN{} crops are the published crop128 protocol. A separate census dropped
\CensusSkipped{} of \CensusTotal{} inputs by output size and left the
\CensusRan{} survivors without reference images to be scored against.

The recommendation is two lines in a results table. Print the trivial arm --- the
input untouched, or the cheapest interpolation --- beside every method, in every
comparison, at every scale. When it loses, nothing is lost by having printed it.
When it wins, print the other endpoint before publishing the ranking, because
what has been learned so far is that the comparison is reporting on its own
measurement, and here the next measurement reversed it.

\section*{Declarations}

%% Funding and competing-interest wording adopted from the owner's packaging
%% script (scripts/build_springer_submission_package.py), shipped in every
%% packet since 2026-08-28 (ruling R16, 2026-09-09). FUNDING_AND_CONFLICTS
%% stays OPEN for the owner's final read.
\subsection*{Funding}

No external funding was received for this study.

\subsection*{Competing interests}

The author declares that they have no competing financial or non-financial
interests that are directly or indirectly related to the work submitted for
publication.

\subsection*{Ethics approval and consent to participate}

Not applicable. This study is a secondary analysis of numeric records produced
by computational runs, together with two predeclared computational runs made
for it: a scoring pass over a fresh sample of a released model on the benchmark
crops, and a toy-scale training ablation. It involves no human participants, no
animal subjects and no personal data, and it reproduces or redistributes no
images. The photographs underlying the benchmark crops are of ordinary scenes
captured for super-resolution research and are distributed under the
benchmark's own terms. Generative restoration models can introduce detail that
was never present in the input, which is a documented hazard wherever a
restored image is used as evidence; the restored outputs produced here were
used only to compute the reported scores, and the paper's central observation
--- that a distortion measure and a perceptual measure can rank the same
outputs in opposite orders --- bears directly on that hazard.

\subsection*{Consent for publication}

Not applicable; no identifiable participant data are included.

\subsection*{Data availability}

The toy training records, the record in which the project withdrew its own toy
comparison, the unpromoted longer run, the per-crop fidelity tables for all
three subset arms, the recorded summaries and stratum counts, the inference
manifests for both attempts, the skip census and its size and source-group
breakdowns, the recorded status of the two unproduced comparisons, the
predeclaration, amendment, per-crop rows, summary, receipt and input provenance
of the perceptual pass, the later VGG secondary scoring pass in a separate
dated directory, and the predeclaration, nine-cell summary, independent
analysis and receipt of the initialisation ablation accompany the manuscript
source. Each artifact is recorded under an evidence ID in a manifest, and the
figure code verifies the bound bytes before reading them, so the build fails
rather than silently reporting stale numbers. The low-resolution inputs and
reference images are from a third-party benchmark held under its own licence
and are cited rather than redistributed; the restored outputs are used only for
scoring and are not included; the model weights and the backbone checkpoint are
third-party releases identified in the manifest by name.
%% Generated by scripts/release_targets.py from the release ledger.
%% Do not edit: an identifier typed here would not be checked against
%% anything, which is exactly the failure the ledger exists to prevent.
%% Ledger state: github=CREATED, zenodo=DEPOSITED
That material is deposited under the persistent identifier \href{https://doi.org/10.5281/zenodo.22647026}{10.5281/zenodo.22647026}, and the development repository is at \url{https://github.com/PeterPonyu/trivial-arm-restoration-audit}.


\subsection*{Code availability}

The figure and table code that emits every number in this manuscript and the
evidence manifest it verifies against are deposited with the data described
above; the two predeclarations, which record the digest of the script each
computation ran, are deposited as bound records. Code is released under the
MIT licence; the manuscript text, figures and recorded result data are released
under CC BY 4.0. The released restoration model and the perceptual-distance
backbone are third-party software, used as distributed and not redistributed.

\subsection*{Materials availability}

Not applicable; the study generated no physical materials.

\subsection*{Author contributions}

\ifdefined\ANONYMOUS
The sole author performed conceptualization, methodology, software, validation,
formal analysis, investigation, visualization, writing---original draft, and
writing---review and editing.
\else
Zeyu Fu: conceptualization, methodology, software, validation, formal analysis,
investigation, visualization, writing---original draft, and writing---review
and editing.
\fi

\ifdefined\DRAFTTODO
\subsection*{Acknowledgements}

\ownertodo{acknowledgements of people, institutions and computing resources
are to be supplied by the author before submission, or this heading deleted;
none is recorded in the project ledger.}
\fi

\subsection*{Use of generative AI tools}

Generative AI tools assisted with code preparation and language editing under
the author's direction. Every number printed above is emitted programmatically
from artifacts tied to the manuscript by cryptographic evidence binding rather
than transcribed into prose, and structural claims --- including the per-crop
means, paired differences, win counts, census, reference-image exclusion, the
anchor rule and the predeclared reading rules --- are enforced by build-time
checks that stop the build. The author reviewed and edited all text and retains
full responsibility for the content.


\bibliographystyle{plain}
\bibliography{references}

\end{document}
");
%% every packet build leaves it undefined, so no unconfirmed declaration text
%% reaches a PDF. The owner replaces each use with the confirmed sentence.
\providecommand{\ownertodo}[1]{\ifdefined\DRAFTTODO[TODO-OWNER: #1]\fi}

%% Self-contained captions. No cross-references, no section numbers: each
%% caption must be readable on its own if the figure is lifted out.

\newcommand{\CapScales}{%
The three scales at which the cheapest arm decided what the comparison could
conclude, placed on one page before those scales are read separately. The
Toy-scale panel restates the toy run: the degraded input sits at
\ToyDegraded{} decibels and the best trained arm at \ToyConv{}, so both
architectures are below doing nothing. The Longer-budget panel restates the
longer budget: the objective falls and fidelity falls with it to \ToyLong{},
still below the untouched input. The Crop-128 panel restates the measured
subset: the released model is behind bicubic on \LossesBicubic{} of
\SubsetN{} crops, by \GapBicubicAbs{} decibels on average, on peak
signal-to-noise ratio. All three panels are fidelity readouts; the perceptual
endpoint on which the subset ranking reverses is drawn on its own figure, and
no panel adds an official ground-truth score.%
}

\newcommand{\CapToy}{%
Three trained arms, all below the input they were asked to repair. A toy-scale
restoration run trained a small convolutional network and a small transformer of
comparable size under a matched step budget on the same data, and scored both
against the same \ToyN{} held-out pairs at \ToyEvalSteps{} sampling steps. Points
are mean peak signal-to-noise ratio of the restored output against the reference
image; higher is closer to the reference. The vertical rule is the same quantity
for the degraded input passed through unchanged, which the same record scores as
an arm in its own right. No trained arm reaches it, so the run supports no
comparison between the two architectures: the shortfall of the best of them is
\ToyConvShortfall{} decibels. The transformer arm's training objective sat at
\ToyTransLoss{} at the last step, against a mean of \ToyTransLastTen{} over the
last ten steps and \ToyConvLoss{} for the convolutional arm at the same step
count: it did not fall, and a denoising objective that stays at unity is what a
network whose output does not depend on its input produces. The stall
reproduces in \InitDefaultStall{} of \InitSeeds{} seeds of the same construction
in a separate predeclared ablation, which also shows that the one
initialisation that zeroes the network's residual gates exactly is the one that
trains; the earlier attribution of the stall to zero-initialised gates is
withdrawn.%
}

\newcommand{\CapBudget}{%
The training budget raised \BudgetFold{}, read on both axes at once. The same
small convolutional restoration network was trained for \ToySteps{} and for
\ToyLongSteps{} steps on the same data at the same learning rate and seed, and
evaluated on the same \ToyN{} held-out pairs at \ToyEvalSteps{} sampling steps.
The two facets carry their own labels: one is the training objective at the last
step, the other the mean peak signal-to-noise ratio of the restored output
against the reference image, where the horizontal rule marks the same quantity
for the degraded input passed through unchanged. The objective falls by a factor
of \ToyLossDrop{}. The fidelity of the output falls with it, by
\ToyLongRegression{} decibels, and stays below the untouched input at both
budgets. Read alone, either facet would give the opposite impression of what the
extra budget bought. The longer run was
written outside the results tree and was never promoted.%
}

\newcommand{\CapSubset}{%
Every crop, the released model against an interpolation. Each point is one crop
of a \SubsetN{}-crop subset drawn from a real-world super-resolution validation
set. The horizontal position is the peak signal-to-noise ratio of bicubic
upsampling of the low-resolution input against the reference image; the vertical
position is the same quantity for a released diffusion restoration model given
the same input, scored by the same metric pass over the same crops. The diagonal
is equality. Points below it are crops on which the interpolation is closer to
the reference. Symbols distinguish the two camera groups the subset is drawn
from, in equal numbers. Peak signal-to-noise ratio measures distance from the
reference and not perceptual quality, so the panel reports which output is
closer to the reference on that measure and not which is better; on a learned
perceptual distance scored on the same crops the order reverses, and that
endpoint is drawn on its own figure. Filled and open symbols identify the two
camera groups; the diagonal is the pre-specified equality reference, not a
fitted trend.%
}

\newcommand{\CapPaired}{%
The paired differences, and the width the model varies by against itself. Panel
A plots one point per crop: the peak signal-to-noise ratio of a released
diffusion restoration model minus that of bicubic upsampling of the same
low-resolution input, over a \SubsetN{}-crop subset, ordered by that difference.
Negative values are crops on which the interpolation is closer to the reference.
The solid rule is zero and the dashed rule is the mean difference. The shaded
band is the largest per-crop shift between two inference passes of the same
model on the same crops, \NoiseMax{} decibels. The record does not show the
sampling seed varied between the passes, so two runs cannot be read as an
estimate of sampling variance; the band bounds only how far re-running the
recorded inference moved a crop. Panel B
gives the mean difference overall and within each camera group, with 95\%
percentile bootstrap intervals over the crops. Camera group is a recorded
property of each crop and was fixed before these differences were computed. The
zero rule is the equality point; the dashed rule is the observed mean, and the
interval in Panel B is a bootstrap summary over this measured subset rather
than uncertainty for the full census.%
}

\newcommand{\CapCeiling}{%
What a memory ceiling selected in an earlier census, and what it removed. That
census ran a released diffusion restoration model over the full set of
\CensusTotal{} inputs
under a fixed budget of \BudgetPixels{} output pixels and a maximum output side
of \BudgetLong{} pixels; inputs whose upsampled output would exceed it were
skipped before inference. Panel A gives, for that census and for the
later \SubsetN{}-crop published crop128 run, what became of every input the run
was given: dropped
before inference, run but left with no reference image to score against, or run
and scored. The census drop and the absence of references were not sampling
decisions; the later scored set is the published protocol, not the remnant of
the skip. Panel B breaks the dropped census inputs down by the recorded source
group
and shades the two groups that appear in the later crop128 set. Every dropped
input
was of a single larger size class, so what survived
the ceiling in the census was fixed by output dimensions rather than chosen.
Segment labels
are n/N and percentages: in Panel A, N is the number of inputs given to that
attempt, and in Panel B it is the total number dropped for memory. An
input that ran without a matched reference is counted as unscored rather than
treated as a zero-fidelity observation; zero-height outcomes remain represented
by the legend.%
}

\newcommand{\CapMetricConsistency}{%
Four fidelity metrics on one paired population. Each point is the mean
DiT--bicubic difference over the same \SubsetN{} crop128 pairs; the vertical
rule is equality and each horizontal whisker is a 95\% percentile bootstrap
interval over those paired differences. The facet labelled PSNR keeps RGB and
luma peak signal-to-noise ratio on their native decibel scale; the facet
labelled SSIM keeps RGB and luma structural similarity on its native unitless
scale. Negative
values mean that the model is farther from the reference under the displayed
fidelity measure. The facets deliberately use separate axes, so their
magnitudes are not compared across units; all four are distortion measures,
and the perceptual endpoint is drawn on its own figure.%
}

\newcommand{\CapCameraSensitivity}{%
Camera-stratified sensitivity of the paired RGB peak signal-to-noise ratio.
Each point is the mean DiT--bicubic difference within one recorded camera
stratum; the horizontal interval is the corresponding descriptive percentile
bootstrap over paired crops, and the label gives the number of DiT wins out of
the stratum's crops. The vertical rule is equality. Both strata are part of the
same 100-crop published crop128 protocol, so this panel does not provide
scene-level or full-census uncertainty and does not measure perceptual quality.%
}

\newcommand{\CapCropCameraMetrics}{%
Crop-level fidelity and scope diagnostics for the published crop128 subset.
Panel A shows the DiT--bicubic difference for every one of the \SubsetN{} paired
crops on RGB PSNR and RGB SSIM, with Canon and Nikon symbols retained. Panel B
gives descriptive 95\% bootstrap intervals for the same two metrics within
each camera stratum; these are intervals over the measured crops, not a
scene-level or full-census uncertainty. Panel C keeps the earlier full-input
census visible: the early \CensusTotal{}-input attempt yielded \CensusRan{} runs
and \CensusSkipped{} skips, whereas the separate crop128 protocol scored
\SubsetN{} crops with references. PSNR and SSIM measure fidelity to a
reference; this figure adds no perceptual endpoint, which is drawn on its own
figure. The metric rows
are recomputed from the bound DiT and bicubic per-crop score files; the bound
camera-stratum record supplies the recorded camera counts and means, and the
census is bound by the earlier evaluation and skip records.}

\newcommand{\CapToyTable}{%
The toy-scale arms as their own training records describe them. Parameter
counts, step counts, wall clock and the training objective at the last step are
read from the manifest each run wrote; fidelity is mean peak signal-to-noise
ratio of the restored output against the reference image over the same
\ToyN{} held-out pairs at \ToyEvalSteps{} sampling steps, and the first row is
the degraded input passed through unchanged, which has no training record.
The two arms trained at \ToySteps{} steps differ by \ToyParamExcess{} per cent in
parameter count, so the budget they share is a step budget and not a capacity
budget, and their wall clocks differ by a factor of \ToySpeedRatio{}, so it is
not a wall-clock budget either.%
}

\newcommand{\CapSubsetTable}{%
The subset comparison, all four recorded fidelity measures. Rows are the two
interpolations and the released diffusion restoration model, each scored over the
same \SubsetN{} crops against the same reference images by the same metric pass.
Peak signal-to-noise ratio and structural similarity are reported over the full
colour image and again over the luma channel with a four-pixel border removed.
The two final rows are paired differences, computed crop by crop and then
averaged, so a negative value means the model was further from the reference than
the interpolation on average. All four measures compare an output with its
reference and none of them measures perceptual quality; the learned perceptual
distance scored on the same crops, on which the order reverses, is reported in
its own figure rather than in this table because it was scored on a fresh
sample of the model rather than on the pixel set behind these rows.%
}

\newcommand{\CapPerceptualReversal}{%
The same crops on both axes of the perception--distortion plane. Each point is
one of the \SubsetN{} crops of the published RealSR crop128 protocol. The
horizontal position is the peak signal-to-noise ratio of a fresh sample of the
released diffusion restoration model, drawn at sampling seed \LpipsSeed{}, minus
that of bicubic upsampling of the same low-resolution input, so negative values
are crops on which the interpolation is closer to the reference; the model
loses on \PsnrLossesLocal{} of \SubsetN{}, and the fresh sample's mean sits
\AnchorDeltaDb{} decibels from the bound fidelity record under a predeclared
anchor rule. The vertical position is the learned perceptual distance (LPIPS,
AlexNet backbone, lower is closer) of bicubic minus that of the model on the
same crop and reference, so positive values are crops on which the model is
perceptually closer; the model wins on \LpipsWins{} of \SubsetN{}, with
\LpipsTies{} ties and an exact sign test of \LpipsSignP{}. The two rules are
equality on each axis and divide the crops into four quadrants, whose counts are
printed: on \QuadLossWin{} crops the interpolation is closer on fidelity and
the model is closer perceptually, and the two endpoints agree on only
\QuadLossLoss{} plus \QuadWinWin{} crops. Filled and open symbols are the two
camera groups. Neither axis is a verdict on the model; together they are the
trade-off, observed.%
}

\newcommand{\CapInitAblation}{%
The toy transformer's training objective under three initialisations at three
seeds each. Every cell is the same network, data, batch, learning rate and
\InitSteps{}-step budget as the bound toy run; only the initialisation and the
seed change. The three facets are the library's default construction, whose
gates are not zero and which is the arm behind the bound record; a standard
transformer scheme without any zeroing step, whose gates are not zero; and that
scheme plus exact zeroing of the adaptive-normalisation gates and final
projection. Each line joins the three points the record holds for one seed: the
mean objective over the first ten steps, the mean over the last ten, and the
last step. The solid rule is unity, the value a denoising objective takes when
the network's output does not depend on its input; the dashed rule is the
predeclared stall threshold of \InitThreshold{} on the last-ten mean. The
default stalls in \InitDefaultStall{} of \InitSeeds{} seeds and the standard
scheme in \InitStandardStall{} of \InitSeeds{}; the zero-gated scheme escapes
in \InitZeroEscape{} of \InitSeeds{}, to a last-ten mean of \InitZeroLoss{}
against \InitStallLoss{} for the cells that stalled. All nine cells restore to
below the degraded input at \IdentityPsnr{} decibels: \InitZeroPsnr{} for the
arm that trained against \InitStallPsnr{} for the arms that did not. The panel
supports a statement about this toy's initialisation and none about the
architecture.%
}

\newcommand{\CapCensusTable}{%
Which inputs a memory ceiling removed, by source group. An earlier attempt over
the full set of \CensusTotal{} inputs skipped every input whose upsampled output
would have exceeded \BudgetPixels{} pixels, which removed \CensusSkipped{} of
them. Groups are the recorded source labels, which correspond to the capture
device. The final column records whether that group appears in the
\SubsetN{}-crop subset measured later. Of the \CensusCameras{} groups the ceiling
touched, \AbsentGroups{} are absent from every measured result in this study.%
}
